\documentclass[12pt]{article}

\usepackage[letterpaper, margin=1in]{geometry}
\usepackage[spanish]{babel}
\usepackage[utf8]{inputenc}
\usepackage[T1]{fontenc}
\usepackage{setspace}
\usepackage{array}
\usepackage{longtable}
\usepackage{hyperref}

\usepackage[
    backend=biber,
    style=ieee
]{biblatex}
\addbibresource{references.bib}

\onehalfspacing

\title{
\textbf{Diseño Metodológico para la Evaluación de Familiaridad Conductual en Agentes Virtuales Inteligentes}
}

\author{
Luis Carlos Quesada Rodríguez\\
Universidad de Costa Rica\\
Escuela de Ciencias de la Computación e Informática
}

\date{\today}

\begin{document}

\maketitle

\section{Definición de Condiciones}

\subsection{Condición Experimental: Condición A}

La condición experimental corresponde a un Agente Virtual Inteligente implementado en Godot, utilizando una representación antropomórfica básica compuesta por un avatar genérico y una voz sintética estándar. Tanto la apariencia visual como la voz del agente fueron diseñadas para minimizar asociaciones explícitas con una persona específica, permitiendo que la familiaridad percibida emerja principalmente a partir del comportamiento conversacional del sistema.

La interacción principal con el agente se realizará mediante una interfaz conversacional basada en texto, complementada por soporte de voz cuando sea posible. Durante la interacción, el agente se presentará como un asistente virtual convencional, manteniendo una conversación relativamente natural con el usuario mientras responde a los distintos temas discutidos durante la sesión experimental.

La característica principal de esta condición corresponde a la incorporación de patrones de familiaridad conductual asociados a una persona conocida por el usuario. Estos patrones serán construidos mediante observación y análisis de muestras conversacionales de la persona modelada, incluyendo elementos como tono comunicativo, expresiones recurrentes, estructura de respuesta, nivel de formalidad, ritmo conversacional y estilo general de interacción.

Adicionalmente, el agente podrá utilizar mecanismos de recuperación contextual asociados al perfil conductual modelado, permitiendo incorporar información relevante relacionada con la persona representada durante la conversación. Esto busca reforzar la percepción de familiaridad y coherencia contextual sin depender exclusivamente de patrones lingüísticos superficiales.

La tarea experimental consistirá en una conversación semi-estructurada de aproximadamente cinco minutos entre el participante y el agente virtual. Durante la interacción, el participante será invitado a hablar sobre su día, actividades recientes o temas cotidianos, permitiendo que la conversación se desarrolle de forma relativamente natural.

Este tipo de interacción fue seleccionado debido a que facilita la aparición de patrones sociales y conversacionales relevantes para la percepción de familiaridad, manteniendo al mismo tiempo un contexto consistente entre participantes.

\subsection{Condición de Control: Condición B}

La condición de control corresponde al mismo Agente Virtual Inteligente utilizando la misma representación visual, voz, interfaz conversacional y tarea experimental. Sin embargo, en esta condición el agente utilizará comportamiento conversacional genérico, sin incorporar patrones de familiaridad conductual ni mecanismos de recuperación contextual asociados a una persona conocida.

De esta manera, la única diferencia significativa entre condiciones corresponde al comportamiento comunicativo del agente y al uso de patrones de familiaridad conductual durante la interacción.

\subsection{Diseño Experimental y Justificación}

El experimento utilizará un diseño intra-sujetos con contrabalanceo parcial. Los participantes serán divididos en dos grupos: el primer grupo interactuará inicialmente con la Condición A y posteriormente con la Condición B, mientras que el segundo grupo seguirá el orden inverso.

Este enfoque permite reducir el efecto de variabilidad individual entre participantes y controlar posibles efectos de orden derivados de la interacción repetida con el sistema.

La comparación entre condiciones permite aislar el efecto de los patrones de familiaridad conductual, manteniendo constantes los elementos visuales, técnicos y contextuales del sistema. De esta manera, el estudio busca evaluar si diferencias en percepción de naturalidad, presencia social, respuesta afectiva y atribución de conocimiento contextual pueden emerger principalmente a partir del comportamiento conversacional del agente.









\section{Matriz de Consistencia Metodológica}

La siguiente matriz presenta la relación entre las preguntas de investigación, los constructos evaluados, los instrumentos seleccionados y las tareas asociadas dentro del protocolo experimental.

{
\small
\renewcommand{\arraystretch}{1.3}
\setlength{\emergencystretch}{3em}

\begin{longtable}{|
>{\raggedright\arraybackslash}p{4.5cm}|
>{\raggedright\arraybackslash}p{3cm}|
>{\raggedright\arraybackslash}p{3cm}|
>{\raggedright\arraybackslash}p{4cm}|}

\hline
\textbf{Pregunta de Investigación} &
\textbf{Variable o Constructo} &
\textbf{Instrumento} &
\textbf{Tarea Asociada} \\
\hline
\endfirsthead

\hline
\textbf{Pregunta de Investigación} &
\textbf{Variable o Constructo} &
\textbf{Instrumento} &
\textbf{Tarea Asociada} \\
\hline
\endhead

¿En qué medida los usuarios son capaces de inferir patrones de familiaridad conductual en un Agente Virtual Inteligente en ausencia de indicaciones explícitas sobre la identidad representada? &
Detección implícita de familiaridad conductual &
Escala Likert personalizada &
Conversación natural con el agente sin revelar la existencia de perfiles de familiaridad conductual. \\
\hline

¿En qué medida la incorporación de patrones de familiaridad conductual en un Agente Virtual Inteligente influye en la percepción de naturalidad y presencia social durante la interacción? &
Naturalidad percibida y presencia social &
Godspeed Questionnaire Series &
Interacción social cotidiana con el agente sobre actividades personales y experiencias recientes. \\
\hline

¿Qué tipo de respuesta afectiva generan los agentes con patrones de familiaridad conductual, en términos de valencia emocional y cercanía percibida, en comparación con agentes sin dichos patrones? &
Valencia emocional y cercanía percibida &
SAM (Self-Assessment Manikin) + ítems Likert complementarios sobre cercanía percibida &
Evaluación emocional posterior a una conversación semi-estructurada con el agente virtual. \\
\hline

¿En qué medida un Agente Virtual Inteligente con patrones de familiaridad conductual es percibido como poseedor de conocimiento contextual previo sobre el usuario? &
Conocimiento contextual percibido &
Ítems adaptados de Poletaykina et al. (2025) \cite{poletaykina_knowing_2025} &
Conversación contextualizada sobre experiencias cotidianas seguida de evaluación de comprensión percibida del agente. \\
\hline

\caption{Matriz de consistencia metodológica del experimento.}
\label{tab:matriz-consistencia}

\end{longtable}
}








\section{Justificación Teórica en HCI}

\subsection{Embodiment, Percepción y Variación Conductual}

El agente virtual incorpora una representación antropomórfica básica mediante un avatar genérico y una voz sintética estándar. Esta decisión se fundamenta en trabajos sobre embodiment y percepción social en agentes virtuales, donde se discute cómo ciertos elementos visuales y comportamentales pueden influir en la interpretación de la interacción por parte del usuario.

Sin embargo, la representación visual del agente fue diseñada deliberadamente para minimizar asociaciones explícitas con individuos reales o características altamente distintivas. El objetivo principal consiste en reducir posibles sesgos derivados de la apariencia visual del agente, permitiendo que la familiaridad percibida emerja principalmente a partir del comportamiento conversacional y no de señales visuales preconcebidas.

En este sentido, la propuesta considera también trabajos relacionados con percepción de avatares y posibles variaciones conductuales inducidas por representaciones virtuales, como las discutidas en la literatura del Efecto Proteus. Aunque estos efectos no constituyen la variable principal del estudio, se reconoce la posibilidad de que la percepción del agente influya en la forma en que el participante conversa, responde emocionalmente o interpreta el comportamiento del sistema durante la interacción.

\subsection{Comportamiento Relacional y Señales Conversacionales}

La condición experimental incorpora patrones de familiaridad conductual basados en una persona conocida por el usuario. Esta decisión se relaciona con trabajos sobre comportamiento relacional en agentes virtuales, particularmente aquellos que exploran cómo ciertos patrones conversacionales y señales sociales pueden influir en la percepción de naturalidad durante la interacción.

En este contexto, la propuesta no busca construir relaciones de largo plazo con el usuario, sino analizar si patrones de comunicación asociados a una persona familiar pueden modificar la percepción social del agente incluso durante interacciones breves.

\subsection{Familiaridad Conductual y Modelos Mentales}

La propuesta se fundamenta en la idea de que los usuarios construyen modelos mentales sobre personas conocidas a partir de experiencias previas de interacción. Estos modelos permiten interpretar comportamiento, anticipar respuestas y reducir incertidumbre durante la comunicación social.

A partir de esta perspectiva, el estudio explora si patrones conductuales asociados a una persona conocida pueden activar parcialmente dichos modelos mentales durante la interacción con un agente virtual. En particular, se busca analizar si esta activación puede emerger principalmente a partir del comportamiento conversacional del agente, sin necesidad de similitud visual explícita.

\subsection{Alineación Conductual y Mimetismo}

Diversos trabajos en interacción humano-agente han estudiado mecanismos de alineación y mimetismo conductual, donde el agente adapta ciertos aspectos de su comportamiento durante la interacción para mejorar la comunicación y percepción social.

Sin embargo, gran parte de estos enfoques se basan en adaptación reactiva al usuario presente durante la conversación. En contraste, esta propuesta explora una aproximación basada en patrones conductuales previamente modelados y asociados a una persona conocida por el usuario. De esta manera, el objetivo no consiste en imitar directamente al participante durante la interacción, sino utilizar patrones previamente aprendidos para inducir percepciones de familiaridad.





\printbibliography
\end{document}